Kurs:Algebraische Kurven (Osnabrück 2012)/Arbeitsblatt 26/latex
Kurs:Algebraische Kurven (Osnabrück 2012)/Arbeitsblatt 26/latex


\setcounter{section}{26}


  \zwischenueberschrift{Aufwärmaufgaben}


\inputaufgabe

{}

{


Man gebe für jedes $n$ ein Beispiel von zwei aus der Schule bekannten ebenen algebraischen Kurven, die sich in genau einem Punkt mit Schnittmultiplizität $n$ schneiden.


}

{} {}


\inputaufgabe

{}

{


Betrachte die durch


\mavergleichskette

{\vergleichskette

{ y
}

{ = }{ 2x^4+3x^2-x+1
}

{  }{
}

{    }{
}

{   }{
}

}
{}{}{}
gegebene Kurve mit dem Punkt


\mavergleichskette

{\vergleichskette

{ P
}

{ = }{ (1,5)
}

{  }{
}

{    }{
}

{   }{
}

}
{}{}{.}
Finde eine Koordinatentransformation derart, dass $P$ zum Punkt $(0,0)$ wird und die Tangente an $P$ zur $x$-Achse.


}

{} {}


\inputaufgabe

{}

{


Es sei eine monomiale ebene Kurve


\mavergleichskette

{\vergleichskette

{ C
}

{ = }{ V \left(  X^d -Y^e  \right)
}

{  }{
}

{    }{
}

{   }{
}

}
{}{}{}
\zusatzklammer {mit $d,e$  teilerfremd} {} {}
gegeben. Berechne die Schnittmultiplizität der Kurve mit einer jeden Geraden $G$ durch den Nullpunkt.


}

{} {}


\inputaufgabegibtloesung

{}

{


Bestimme die Schnittmultiplizität im Nullpunkt des Kartesischen Blattes


\mavergleichskettedisp

{\vergleichskette

{ C
}

{ =} { V \left(  X^3+Y^3-3XY  \right)
}

{ } {
}

{ } {
}

{ } {
}

}
{}{}{}
mit jeder affinen Geraden der affinen Ebene. Man setze voraus, dass die Charakteristik des Körpers nicht $3$ ist.


}

{Man setze voraus, dass die Charakteristik des Körpers nicht $3$ ist.} {}


  \zwischenueberschrift{Aufgaben zum Abgeben}


\inputaufgabe

{}

{


Berechne die
\definitionsverweis {Schnittmultiplizität}{}{}
der beiden
\definitionsverweis {monomialen Kurven}{}{}


\mathdisp {C=V \left(  X^5 -Y^2  \right) \text{ und } D=V \left(  X^7 -Y^3  \right)} {  }


im Nullpunkt.


}

{} {}


\inputaufgabe

{}

{


Es sei $K$ ein
\definitionsverweis {Körper}{}{}
und seien


\mavergleichskette

{\vergleichskette

{ C
}

{ = }{ V(F)
}

{  }{
}

{    }{
}

{   }{
}

}
{}{}{}
und


\mavergleichskette

{\vergleichskette

{ D
}

{ = }{ V(G)
}

{  }{
}

{    }{
}

{   }{
}

}
{}{}{}
\definitionsverweis {ebene algebraische Kurven}{}{.}
Es sei


\mavergleichskette

{\vergleichskette

{ P
}

{ \in }{ C
}

{  }{
}

{    }{
}

{   }{
}

}
{}{}{}
ein glatter Punkt, sodass der
\definitionsverweis {lokale Ring}{}{}


\mavergleichskette

{\vergleichskette

{ R
}

{ = }{ (K[X,Y]_{ {\mathfrak m} } )/(F)
}

{  }{
}

{    }{
}

{   }{
}

}
{}{}{}
ein
\definitionsverweis {diskreter Bewertungsring}{}{}
ist. Zeige, dass die Beziehung


\mavergleichskettedisp

{\vergleichskette

{ \operatorname{mult} _{ { P} }  (   F ,  G  )
}

{ =} { \operatorname{ord} \, (G)
}

{ } {
}

{ } {
}

{ } {
}

}
{}{}{}
gilt, wobei $\operatorname{ord} \,$ die Ordnung im Bewertungsring $R$ bezeichnet.


}

{} {}


\inputaufgabe

{}

{


Betrachte die Parabel


\mavergleichskette

{\vergleichskette

{ C
}

{ = }{ V \left(  Y-X^2  \right)
}

{  }{
}

{    }{
}

{   }{
}

}
{}{}{}
und den Kreis $D$ mit Mittelpunkt $(0,r)$ und Radius $r$. Bestimme die Schnittpunkte von $C$ und $D$ und die jeweiligen Schnittmultiplizitäten.


}

{} {}


\inputaufgabe

{}

{


Bestimme für den
\definitionsverweis {Restklassenring}{}{}


\mathl{{\mathbb C}[X,Y]/ \left(  XY-1,X^2+Y^2-a  \right)}{}
\zusatzklammer {für jedes


\mavergleichskettek

{\vergleichskettek

{ a
}

{ \in }{ {\mathbb C}
}

{  }{
}

{    }{
}

{   }{
}

}
{}{}{}} {} {}
eine Beschreibung als Produktring von lokalen Ringen. Man gebe dabei die ${\mathbb C}$-Dimensionen der beteiligten Ringe an.


}

{} {}


\inputaufgabe

{}

{


Bestimme für die
\definitionsverweis {Kurve}{}{}


\mathl{V \left(  X^3+Y^3-3XY+1  \right)}{} die
\definitionsverweis {singulären Punkte}{}{}
über $\R$ und über ${\mathbb C}$. Man gebe jeweils die
\definitionsverweis {Multiplizität}{}{}
und die
\definitionsverweis {Tangenten}{}{}
an.


}

{} {}


\inputaufgabe

{}

{


Betrachte die durch


\mavergleichskette

{\vergleichskette

{ y
}

{ = }{ 2x^4+3x^2-x+1
}

{  }{
}

{    }{
}

{   }{
}

}
{}{}{}
gegebene Kurve im Punkt


\mavergleichskette

{\vergleichskette

{ P
}

{ = }{ (1,5)
}

{  }{
}

{    }{
}

{   }{
}

}
{}{}{}
in den in
Aufgabe 26.2
gefundenen Koordinaten. Bestimme die Potenzreihe für die Kurve in $P$ entlang der Tangente.


}

{} {}


Die folgende Aufgabe ist vermutlich schwieriger.


\inputaufgabe

{}

{


Es seien zwei verschiedene monomiale ebene Kurven


\mavergleichskette

{\vergleichskette

{ C
}

{ = }{ V \left(  X^d -Y^e   \right)
}

{  }{
}

{    }{
}

{   }{
}

}
{}{}{}
und


\mavergleichskette

{\vergleichskette

{ D
}

{ = }{  \left(  X^r -Y^s  \right)
}

{  }{
}

{    }{
}

{   }{
}

}
{}{}{}
gegeben (mit $d,e$ und $r,s$ teilerfremd). Berechne die Schnittmultiplizität der beiden Kurven im Nullpunkt.


}

{} {}


<< | Kurs:Algebraische Kurven (Osnabrück 2012) | >>


PDF-Version dieses Arbeitsblattes


 Zur Vorlesung
(PDF)
